% ============================================================
% PROGRAMMATION JULIA
% Résumé de cours (version française) — 30 heures
% Compilation : xelatex -shell-escape cours.tex (2 passes)
%
% Ce fichier est une version française abrégée : préface,
% table des matières, et aperçu chapitre par chapitre. Le
% contenu détaillé (code, preuves, exercices) est pour
% l'instant disponible uniquement dans la version anglaise
% notes.pdf au dossier parent en/. Une traduction complète
% des 10 chapitres est prévue ultérieurement.
% ============================================================
\documentclass[12pt,a4paper]{article}

\usepackage{fontspec}
\usepackage{polyglossia}
\setmainlanguage{french}

\usepackage[margin=2.5cm]{geometry}
\usepackage{hyperref}
\hypersetup{
  colorlinks=true,
  linkcolor=blue!60!black,
  urlcolor=blue!60!black,
}
\usepackage{enumitem}

\title{\textbf{Programmation Julia}\\[0.4em]
       \large Résumé de cours \textendash{} 30 heures}
\author{Yaé Ulrich Gaba}
\date{}

\begin{document}
\maketitle

\section*{Préface}

Ce document est la version française abrégée du cours \emph{Julia Programming}. La version anglaise complète \textendash{} dix chapitres, code exécutable, exercices et projets \textendash{} se trouve dans le fichier \texttt{../en/notes.pdf} du présent dépôt. La traduction française intégrale est prévue mais pas encore disponible.

Le cours couvre Julia depuis les fondamentaux du langage jusqu'au calcul scientifique et à l'apprentissage automatique~: types, dispatch multiple, tableaux et algèbre linéaire, manipulation de données avec DataFrames, visualisation, optimisation de performance, calcul scientifique (équations différentielles via \texttt{DifferentialEquations.jl}), et apprentissage automatique avec \texttt{MLJ.jl} et \texttt{Flux.jl}.

Public visé~: étudiants de licence ou de master en mathématiques, physique, ou science des données ayant une expérience préalable d'un langage comme Python, R, MATLAB ou C++. Prérequis mathématiques~: algèbre linéaire et calcul différentiel au niveau licence.

\section*{Pourquoi Julia~?}

Julia a été conçu pour résoudre le \emph{problème des deux langages}~: les chercheurs prototypent en langage de haut niveau (Python, MATLAB, R) puis réécrivent le code critique en termes de performance en C ou Fortran. Julia élimine cette étape~: le même code, écrit dans un style proche de Python, atteint des vitesses comparables à C grâce à la compilation JIT via LLVM.

Trois caractéristiques définissent Julia~:
\begin{itemize}[leftmargin=1.5em]
  \item \textbf{Compilation JIT.} Les fonctions sont compilées en code machine natif lors de leur premier appel. Les appels suivants sont rapides.
  \item \textbf{Dispatch multiple.} Les fonctions sont spécialisées pour chaque combinaison de types d'arguments, ce qui remplace la programmation orientée objet par un modèle plus proche de la pensée mathématique.
  \item \textbf{Langage conçu pour le calcul scientifique.} Syntaxe mathématique naturelle, Unicode pour les variables (\verb|α|, \verb|β|, \verb|∇|), bibliothèques numériques et scientifiques de première classe.
\end{itemize}

\section*{Aperçu des chapitres}

\begin{enumerate}[leftmargin=2em, itemsep=0.4em]
  \item \textbf{Prise en main de Julia.} Installation, JuliaUp, REPL, gestionnaire de paquets \texttt{Pkg}, VS Code, Pluto, et premier programme exécutable.
  \item \textbf{Types, variables et structures de contrôle.} Système de types, hiérarchie des nombres, structures composites, conditionnelles et boucles.
  \item \textbf{Fonctions et dispatch multiple.} Définition de fonctions, méthodes, dispatch sur les types d'arguments, et son rôle central dans l'écosystème Julia.
  \item \textbf{Tableaux et algèbre linéaire.} Tableaux multidimensionnels, indexation, opérations vectorisées, bibliothèque \texttt{LinearAlgebra} standard.
  \item \textbf{DataFrames et manipulation de données.} \texttt{DataFrames.jl}, chargement de fichiers CSV, filtres, jointures, opérations \texttt{group-by}.
  \item \textbf{Visualisation.} \texttt{Plots.jl} comme couche d'abstraction, \texttt{Makie.jl} pour les rendus de qualité publication, choix d'un back-end selon le contexte.
  \item \textbf{Performance et optimisation.} Annotations de types, pièges classiques (variables globales, instabilité de type), profilage, benchmarking avec \texttt{BenchmarkTools.jl}.
  \item \textbf{Calcul scientifique.} Équations différentielles ordinaires et stochastiques via \texttt{DifferentialEquations.jl}, optimisation, calculs symboliques avec \texttt{Symbolics.jl}.
  \item \textbf{Apprentissage automatique avec Julia.} Écosystème \texttt{MLJ.jl} (modèles classiques), \texttt{Flux.jl} (apprentissage profond), différentiation automatique via \texttt{Zygote.jl}.
  \item \textbf{Projets capstone.} Trois projets intégrés~: pricing d'options avec Monte-Carlo, classification d'images MNIST, modélisation épidémiologique de la propagation d'une maladie infectieuse.
\end{enumerate}

\section*{Pour aller plus loin}

\begin{itemize}[leftmargin=1.5em]
  \item \textbf{Version anglaise complète}~: \texttt{../en/notes.pdf} (environ 200 pages avec code et exercices).
  \item \textbf{Documentation officielle Julia}~: \url{https://docs.julialang.org}.
  \item \textbf{Forum communautaire}~: \url{https://discourse.julialang.org}.
  \item \textbf{Cours gratuit du MIT}~: \emph{Introduction to Computational Thinking}, \url{https://computationalthinking.mit.edu}.
\end{itemize}

\section*{Contact}

Pour toute question, suggestion ou demande concernant la traduction complète, écrire à~: \href{mailto:gabayae2@gmail.com}{\texttt{gabayae2@gmail.com}}.

\end{document}
